\SourceLeaf{307}
\Needspace{8\baselineskip}\begingroup\small
\begin{center}
BEMERKUNGEN.
\end{center}

Auf die in dem Briefe vom 12. Mai 1843 erwähnten Untersuchungen bezieht sich die Notiz: \emph{Zur Theorie der Kegelschnitte}, die aus dem Nachlass in der folgenden Abtheilung: Verwendung complexer Grössen für die Geometrie S. 341 bis 344 dieses Bandes abgedruckt ist.

\textsc{Gauss}' Lösung der von \textsc{Möbius} gestellten geometrischen Aufgabe hat \textsc{Schumacher} in den Astronomischen Nachrichten No. 42, 1823 November veröffentlicht; sie ist wiederabgedruckt in Bd. IV. S. 406---407.

Die Notiz [II] steht auf einem einzelnen Blatte; sie stammt wahrscheinlich aus derselben Zeit, wie die unter [I] abgedruckten Briefstellen. \textsc{Gauss} hat den barycentrischen Calcul, wie unter [3] angedeutet, insbesondere auf die Theorie des Fünfecks angewandt. Die betreffenden Entwicklungen, die mit dem Pentagramma mirificum in Zusammenhang stehen, findet man unter dem Titel: \emph{Weitere Fragmente über das Pentagramma mirificum} No. 11, S. 109 bis 111 dieses Bandes (vergl. auch die zugehörigen Bemerkungen \textsc{Frickes}, S. 117). Auch auf die Notiz: \emph{Beziehung der Raumverhältnisse auf ein gegebenes Tetraeder}, Bd. II. S. 307 bis 308 möge hingewiesen werden; sie stammt wahrscheinlich schon aus dem Juli 1831.

\begin{flushright}
\textsc{Stäckel}.
\end{flushright}
\par\endgroup
